\chapter{Introduction}
\label{ch:introduction}

\section{Background}
\label{sec:background}

Chronic kidney disease (CKD) is a progressive, largely irreversible loss of
kidney function that affects an estimated 800 million people worldwide, or
roughly one adult in ten \cite{cdcbasics,kdigo2024}. It is a leading cause of
cardiovascular mortality, and in its terminal stage it requires dialysis or
transplantation---interventions that are financially and logistically out of
reach for a large proportion of the population in low- and middle-income
countries such as Bangladesh.

The clinical tragedy of CKD is that it is both \emph{slow} and \emph{silent}.
The disease typically develops over years, and a person can lose a substantial
fraction of kidney function before experiencing any symptom that would prompt a
medical visit. By the time symptoms such as oedema, persistent fatigue, or
reduced urine output appear, the disease has often reached a stage at which the
remaining clinical options are limited to slowing further decline rather than
restoring function. Early identification is therefore the single most effective
lever available, because the dominant drivers of CKD progression---elevated
blood pressure, poorly controlled blood glucose, and central obesity---are
precisely the factors that respond to early lifestyle modification and routine
clinical follow-up.

\subsection{How CKD is confirmed clinically}

The internationally accepted framework for staging CKD is the KDIGO risk grid
\cite{kdigo2024}, which combines two laboratory measurements:

\begin{itemize}
  \item the \textbf{estimated glomerular filtration rate} (\eGFR), computed
        from serum creatinine, which measures how effectively the kidneys
        filter blood; and
  \item the \textbf{urine albumin-to-creatinine ratio} (\UACR), which measures
        how much protein is leaking into the urine.
\end{itemize}

Together these two axes define the low, moderate, high and very-high risk
categories shown later in Figure~\ref{fig:kdigo}. This grid is the diagnostic
gold standard, and nothing in this thesis proposes to replace it.

The practical difficulty is one of \emph{ordering}, not of accuracy. Neither
\eGFR{} nor \UACR{} is part of a routine general health checkup. A creatinine
assay with an \eGFR{} calculation, and particularly a urine albumin test, are
kidney-specific orders that a clinician places when they already have a reason
to suspect kidney disease. For an asymptomatic patient there is usually no such
trigger. The result is a circular detection gap: the test that would reveal
early CKD is not ordered because there is no symptom, and there is no symptom
because the disease is still early. This gap---not the cost of a blood draw as
such---is the specific problem this research addresses.

\subsection{What a routine checkup already provides}

A standard general health checkup, by contrast, routinely produces a
substantial amount of information that is physiologically related to kidney
health: blood pressure, waist circumference, fasting glucose, glycated
haemoglobin (HbA1c), blood urea nitrogen (BUN), uric acid, serum albumin and
haemoglobin, together with a self-reported history of diabetes and
hypertension. These measurements are inexpensive, widely available, and---
critically---\emph{already being collected} in general practice, community
health camps and population health surveys. If a sufficiently reliable risk
signal can be extracted from measurements that already exist, then screening
becomes an act of computation rather than an act of additional testing.

\subsection{Language models as a patient-education layer}

A risk score on its own has limited value to the person being screened. Being
told ``your kidney risk is moderate'' answers neither ``what does that mean for
me?'' nor ``what should I do next?''. Health literacy about kidney disease is
low in the general population, and the gap between a numeric model output and
an actionable understanding is precisely where most screening tools stop.

Recent instruction-following large language models (LLMs) can produce fluent,
readable explanations, and are therefore an attractive candidate for closing
that gap. They are also, however, unreliable in exactly the ways that matter
most in a medical context: they can state a diagnosis that the underlying model
never made, recommend or adjust medication, invent numeric dietary targets, or
reassure a patient whose reported vital signs indicate an emergency. Deploying
an LLM in a health-screening pathway is therefore not primarily a generation
problem but a \emph{constraint} problem.

\section{Overview}
\label{sec:overview}

This thesis presents \textbf{Nephro-AI}, an end-to-end system that pairs a
cross-sectional machine-learning kidney-risk screening model with a
safety-constrained, guideline-grounded language-model education layer. The
system is built around two design commitments that run through every chapter
that follows.

\begin{enumerate}
  \item \textbf{The screening model consumes only routine, low-cost inputs.}
        Creatinine, \eGFR, \UACR{} and urine protein are excluded from the
        feature set entirely. They are used exclusively to construct the
        training label from the KDIGO grid, and never as model inputs. The
        deployed model therefore runs on data a general checkup already
        produces.
  \item \textbf{Safety is enforced deterministically, not merely requested.}
        The language model is instructed to behave safely, is fine-tuned on
        safety-screened examples, and---decisively---has every one of its
        outputs passed through a rule-based validator that blocks diagnosis,
        medication and supplement advice, rewrites numeric dietary targets, and
        injects urgent-care redirection when the reported vitals demand it.
\end{enumerate}

The complete system is implemented as a working web application with a
patient-facing interface and a researcher-facing audit interface, backed by a
FastAPI service, a PostgreSQL research database, and a local language-model
runtime. Both the screening model and the education layer are evaluated
quantitatively, the latter through a pre-registered six-arm ablation designed
to attribute the observed safety behaviour to specific architectural
components rather than to the system as an undifferentiated whole.

\section{Problem Statement}
\label{sec:problem}

The problem addressed by this research can be stated as follows.

\begin{quote}
\emph{Given only the routine, low-cost measurements produced by a standard
general health checkup---and explicitly excluding the kidney-specific
confirmatory tests (\eGFR, creatinine, \UACR)---can a machine-learning system
stratify an adult into a clinically meaningful kidney-risk category with
sufficient reliability to justify referral for confirmatory testing; and can
the resulting risk assessment be communicated to that person as personalised,
guideline-grounded education by a large language model whose output is
guaranteed, by construction, never to diagnose, prescribe, or withhold urgent
care?}
\end{quote}

Two sub-problems follow. The first is a supervised learning problem under
severe class imbalance: in the population data used here, 82.9\% of
participants fall in the low-risk category and only 4.5\% in the high-risk
category, so a classifier optimised for aggregate accuracy will systematically
under-detect exactly the group that most needs detection. The second is a
constrained natural-language generation problem: the education text must be
simultaneously personalised (referring to the individual's actual measurements),
grounded (traceable to published clinical guidance), readable (approximately
eighth-grade level), and unconditionally safe.

\section{Motivation}
\label{sec:motivation}

Three observations motivated this work.

\textbf{First, the existing machine-learning literature on CKD is
methodologically circular for screening purposes.} A large body of published
work reports very high accuracy---frequently above 95\%---for CKD
classification \cite{chittora2021,aljaaf2019,yildirim2017}. On inspection,
these models almost universally include serum creatinine, blood urea, or
\eGFR{} among their input features, and are frequently evaluated on a small,
clean, well-balanced benchmark dataset. A model that consumes creatinine to
predict a creatinine-derived label is performing a computation, not a
screening: it requires the very test whose ordering is the bottleneck. The
reported accuracy is real but the clinical utility for early detection is
limited.

\textbf{Second, screening tools stop at the score.} Systematic review of kidney
prediction models \cite{ramspek2021} finds that the overwhelming majority
terminate at a risk estimate, with no patient-facing communication component
and no evaluation of whether the information conveyed is understood or
actionable. The educational half of the problem is treated as out of scope.

\textbf{Third, health-domain LLM systems rarely publish measured safety
guarantees.} Where language models are deployed for patient-facing health
information, safety is typically pursued through prompt engineering alone,
which offers no guarantee, or through fine-tuning, whose effect is difficult to
verify. Published, component-attributed, statistically tested measurements of
unsafe output rates are uncommon. Since safety is the property on which
deployment permission ultimately depends, it deserves the same evaluative rigour
normally reserved for accuracy.

\section{Objectives}
\label{sec:objectives}

The overall goal of this research is to design, implement and evaluate a
kidney-risk screening and education system that is deployable on routinely
available data and demonstrably safe in its patient-facing output. The specific
objectives are:

\begin{enumerate}
  \item To construct a KDIGO-aligned, cross-sectional kidney-risk dataset from
        NHANES 2021--2023 in which the target label is derived from \eGFR{} and
        \UACR{} while the feature set contains neither.
  \item To design and train a two-stage classification cascade that is
        explicitly optimised to minimise the clinically most dangerous error---
        classifying a genuinely high-risk person as low risk---rather than to
        maximise aggregate accuracy.
  \item To produce a per-patient explanation of each prediction and a
        transparent reliability indicator, so that the screening output is
        interpretable rather than opaque.
  \item To construct a safety-scoped knowledge base from published clinical
        guidance (KDIGO, ADA, NKF/KDOQI, CDC) and a retrieval mechanism that
        selects guidance according to the individual's own dominant risk
        factors.
  \item To adapt a compact open-weight language model, by parameter-efficient
        fine-tuning, to produce structured, personalised, guideline-grounded
        education from a short prompt.
  \item To implement a deterministic post-generation safety validator that
        enforces the clinical constraints unconditionally, independently of
        what the language model produces.
  \item To evaluate the education layer with a pre-registered ablation that
        attributes the measured safety behaviour to individual components, using
        appropriate paired statistical inference.
  \item To deliver the whole system as a working, auditable web application.
\end{enumerate}

This work also seeks to answer the following research questions:

\begin{enumerate}
  \item How much kidney-risk discrimination is recoverable from routine, low-cost
        inputs alone, once the confirmatory markers are withheld?
  \item Does decomposing the three-class problem into a cascade of two binary
        decisions reduce the rate of high-to-low misclassification relative to a
        single multi-class classifier?
  \item Which architectural component---retrieval, fine-tuning, or the
        deterministic validator---is actually responsible for the safety of the
        delivered patient education?
  \item Can a fine-tuned model reproduce full-prompt safety behaviour from a
        substantially shorter prompt, and at what measurable cost?
\end{enumerate}

\section{Scope, Assumptions and Limitations}
\label{sec:scope}

The following assumptions delimit the work.

\begin{itemize}
  \item \textbf{Screening, not diagnosis.} The system outputs a risk category
        intended to trigger confirmatory testing. It does not diagnose CKD, and
        every patient-facing surface states this explicitly.
  \item \textbf{Cross-sectional design.} Both the training data and each
        individual screening are point-in-time. NHANES provides a single
        observation per participant, so the model cannot represent trajectory,
        and the system does not track change across visits.
  \item \textbf{Label fidelity.} The target variable is a KDIGO-grid category
        computed from measured \eGFR{} and \UACR{} within the survey. It is a
        laboratory-derived risk stratification, not an adjudicated clinical
        diagnosis with imaging or biopsy confirmation.
  \item \textbf{Population.} The model is trained on a United States adult
        survey population. Transfer to a Bangladeshi or other South Asian
        population is untested and would require local validation.
  \item \textbf{Language-model evaluation is synthetic.} The 300-profile
        benchmark used to evaluate the education layer is generated and scored
        against the deployed cascade, not drawn from real patients, and no
        clinician review of the generated text was conducted in this phase.
  \item \textbf{Rule-based safety detection.} Both the deployed validator and
        the evaluation detectors are deterministic rule systems. They provide
        reproducibility and auditability, but cannot be assumed complete against
        an unbounded space of unsafe phrasings.
\end{itemize}

\section{Contributions}
\label{sec:contributions}

The principal contributions of this thesis are:

\begin{enumerate}
  \item \textbf{A routine-input kidney-risk cascade.} A two-stage
        random-forest/gradient-boosting cascade trained on 4{,}719 NHANES
        participants using fourteen routine inputs, achieving 70.9\% accuracy
        and 73.7\% high-risk recall on a held-out test set of 833 participants
        with only two high-to-low misclassifications, without access to
        creatinine, \eGFR{} or \UACR.
  \item \textbf{Factor-aware retrieval for per-patient grounding.} A retrieval
        strategy that derives a priority topic ordering from the individual's own
        model-attributed risk factors and guarantees topic coverage, so that two
        people in the same risk category receive different guideline grounding.
  \item \textbf{A prompt-compression fine-tune.} A LoRA adapter trained on
        540 safety-screened examples in which targets generated under a long
        instruction prompt are paired with a short prompt, reducing prompt cost
        by 25\% at serving time without a measurable safety penalty.
  \item \textbf{A component-attributed safety evaluation.} A pre-registered
        six-arm ablation over a frozen 300-profile benchmark, analysed with
        McNemar tests, Wilson intervals and Holm correction, which isolates the
        deterministic validator---rather than fine-tuning---as the component
        responsible for eliminating critical-unsafe output.
  \item \textbf{A complete, auditable implementation}, including a
        patient-facing interface, a researcher-facing audit interface exposing
        retrieval and validator decisions, and a reproducible artefact trail.
\end{enumerate}

\section{Research Outline}
\label{sec:outline}

The remainder of this report is structured as follows.
\textbf{Chapter~\ref{ch:literature}} reviews the clinical background of CKD
staging, the machine-learning literature on CKD prediction, techniques for
learning under class imbalance, model explainability, and the use of large
language models, retrieval augmentation, parameter-efficient fine-tuning and
safety guardrails in healthcare, closing with a statement of the research gap.
\textbf{Chapter~\ref{ch:system}} presents the system model: the layered
architecture, the formal definition of the cascade, the explanation and
reliability formulations, and the end-to-end request algorithm.
\textbf{Chapter~\ref{ch:methodology}} details the methodology---dataset
construction and label derivation, feature selection, imbalance handling, model
training and threshold selection, knowledge-base construction, retrieval,
prompt engineering, training-data generation, fine-tuning, and the safety
validator. \textbf{Chapter~\ref{ch:implementation}} describes the
implementation: the technology stack with its trade-offs, the API surface, the
web interfaces, the database schema and the deployment modes.
\textbf{Chapter~\ref{ch:results}} reports results for both components,
including the six-arm ablation and its statistical analysis.
\textbf{Chapter~\ref{ch:discussion}} interprets these findings, examines
threats to validity, and discusses ethical and deployment considerations.
\textbf{Chapter~\ref{ch:conclusion}} concludes and identifies future work.
